\section{A stationary sample of hard-core paths}\label{path:section}

The spatial estimates will supply many particles whose short trajectories
can later be changed.  We first represent the position law by an exact
killed-Brownian sample, and prove that prescribed cells simultaneously lack
the required trajectories with very small probability.  All probabilities
in this section refer either to that sample or to explicitly specified free
Brownian paths; the interacting trajectories are never assumed independent.

We use the scaled torus, cells, and parameters of
Equation~\eqref{setup:scales}.  In particular,
\[
 V=K^3,\qquad D=N/V,\qquad r=\alpha/D,
 \qquad \alpha\in[\alpha_0/2,2\alpha_0],\qquad
 \mu(\dd X)=\Phi(X)^2\dd X.
\]
The constants $\alpha_0,\ell,\theta,b$ have the meanings fixed in the
spatial estimates, and $s_0=bD^{-1/3}$.  Choose $\theta$ with
$3\theta<\ell/2$ before fixing $b$ as in
Proposition~\ref{sp:joint-tails}.  The time $t\in(0,1]$ will be small but
fixed before the lower threshold on $D$ is chosen.

\subsection{The endpoint-weighted law}

For the killed Dirichlet semigroup on arbitrary open Euclidean domains,
see \cite[arXiv version~2, Proposition~2.9 and Theorem~6.1]{BroderixHundertmarkLeschke2000Semigroups}.
That reference uses Brownian generator $\Delta/2$; our convention is
$\Delta$. The exhaustion argument below gives the required torus realization.

Let $\nu_N$ be the measure on $N$ independent torus Brownian motions
of generator $\Delta$, on $[0,2t]$, with Lebesgue initial measure.
Write $\omega=(\omega_1,\ldots,\omega_N)$ and
$X(u)=(\omega_1(u),\ldots,\omega_N(u))$.  Set $Y=X(0)$ and $Z=X(2t)$.
Define
\begin{equation}\label{path:law}
 \dd\cR(\omega)=e^{2tE_N}\Phi(Y)\Phi(Z)
 \ind_{\{d(\omega_i(u),\omega_j(u))>r\ \forall u,\ i<j\}}
 \dd\nu_N(\omega).
\end{equation}
Here and below $d$ is the torus distance.  A continuous path has a
continuous Euclidean lift after its initial representative is chosen;
its lifted increments and diameter do not depend on that representative.

\begin{lemma}[Stationary killed paths]\label{path:slab}
Equation~\eqref{path:law} defines a probability measure.  Every time
marginal is $\mu$.  Its restriction to an interval of length $s$ has
the same construction with endpoint multiplier
$e^{sE_N}\Phi(X_-)\Phi(X_+)$, and it is invariant under time reversal.
Given both endpoint configurations, it is the product torus-bridge law
conditioned on mutual hard avoidance.  If the winding increments are
also fixed, the underlying product consists of the corresponding
Euclidean bridges.
\end{lemma}

\begin{proof}
For the Dirichlet Laplacian on the open allowed configuration set
$\Omega$, the killed heat kernel represents $e^{-sH_N}$.
This representation applies to the exact form domain $H^1_0(\Omega)$.
One way to see that no regularity assumption on the collision boundary
is needed is to exhaust $\Omega$ by increasing smooth relatively
compact open sets.  On each set the stopped heat equation gives the
killed representation.  The union of their form domains is dense in
$H^1_0(\Omega)$, since every smooth compactly supported function is
supported in one exhaustion member.  The variational characterization
of resolvents therefore gives strong resolvent convergence, and then
strong heat-semigroup convergence.  A continuous path remaining in
$\Omega$ throughout a compact time interval has compact image in
$\Omega$, hence lies in one exhaustion member.  Monotone convergence
on paths identifies the limiting kernel.

The integral in Equation~\eqref{path:law} is
$e^{2tE_N}\langle\Phi,e^{-2tH_N}\Phi\rangle=1$.
Splitting the killed kernel at time $u$ gives the marginal density
\[
 e^{2tE_N}(e^{-uH_N}\Phi)(x)
                  (e^{-(2t-u)H_N}\Phi)(x)=\Phi(x)^2.
\]
Integrating an added initial or terminal interval proves the restriction
statement.  Symmetry of the killed kernel proves reversal.  At fixed
endpoints the two factors $\Phi$ are constant, so disintegration of
the free measure leaves exactly its independent bridges and the
avoidance indicator.  Resolving the endpoint heat kernel into its
Euclidean translates gives the last assertion.  All conditional
statements are understood for almost every datum under its stated law.
\end{proof}

The next inequality transfers a rare event for one free trajectory to
an estimate for any specified set of labelled trajectories.  It uses
only the removal-energy estimate, not independence under $\cR$.

\begin{lemma}[Dropping selected interactions]\label{path:deletion}
Let $I\subset\{1,\ldots,N\}$ have cardinality $m$.  If $h$ is a
nonnegative measurable single-path functional invariant under reversal,
then
\begin{equation}\label{path:deletion-bound}
 \E_{\cR}\prod_{i\in I}h(\omega_i)
 \le e^{Ct\alpha m}
       \int\mu(\dd Y)\prod_{i\in I}\E^{Y_i}_{\rm free}h.
\end{equation}
The free expectation uses a Brownian path of length $2t$.
\end{lemma}

\begin{proof}
First take bounded $h$.  Drop every collision restriction involving a
label in $I$.  For fixed selected endpoint vectors $y,z$, integrate
the other paths with their killed semigroup.  The resulting matrix
element between the corresponding slices of $\Phi$ is at most
\[
 e^{-2tE_{N-m}}a(y)a(z),\qquad
 a(y)^2=\int\Phi(y,x')^2\dd x'.
\]
The free selected paths, weighted by $\prod h$, give a symmetric
endpoint kernel, by reversal invariance.  More explicitly, let $\nu_I$
be their product free-path measure with Lebesgue initial measure, and
write $y=\omega_I(0)$ and $z=\omega_I(2t)$.  Then
\begin{align*}
 \int a(y)a(z)\prod_{i\in I}h(\omega_i)\dd\nu_I
 &\le\frac12\int\bigl(a(y)^2+a(z)^2\bigr)
                       \prod_{i\in I}h(\omega_i)\dd\nu_I\\
 &=\int a(y)^2\prod_{i\in I}\E_{\rm free}^{y_i}h\,\dd y.
\end{align*}
Here the equality reverses the paths in the term containing $a(z)^2$;
independence is used only under the free measure.  By the definition of
$a$, the last integral is
$\int\mu(\dd Y)\prod_{i\in I}\E_{\rm free}^{Y_i}h$.
Reinstating the factor $e^{2t(E_N-E_{N-m})}$ and applying
Corollary~\ref{sp:removal} proves Equation~\eqref{path:deletion-bound}.
Applying this to $h\wedge M$
and letting $M$ increase proves the general case.
\end{proof}

\subsection{Uniform path cutoffs}

Put
\begin{equation}\label{path:cutoff-parameters}
 \gamma=\frac1{12},\qquad h_0=D^{-1},\qquad
 m_t=\lceil t/r^2\rceil,\qquad \tau=t/m_t.
\end{equation}
We henceforth take $D$ large enough that $h_0<t$ and
$r^2/2\le\tau\le r^2$.  The two halves of $[0,2t]$ are divided into
$m_t$ intervals each.  If these intervals are $I_1,\ldots,I_{2m_t}$,
write
\[
 A_j(\omega)=r^{-1}\diam(\widetilde\omega(I_j)),
\]
where $\widetilde\omega$ is a continuous lift.

\begin{definition}\label{path:regular-path}
Fix a sufficiently large numerical constant $C_{\rm reg}$.
A path is \emph{regular} if
\begin{equation}\label{path:regular-conditions}
 \max_j A_j\le D^\gamma,
 \qquad \frac1{2m_t}\sum_{j=1}^{2m_t}(1+A_j)^3\le C_{\rm reg},
\end{equation}
and its lifted diameter on each of $[0,h_0]$ and
$[2t-h_0,2t]$ is at most $s_0/8$.
\end{definition}

For two regular paths whose corresponding outer endpoints are more than
$s_0$ apart, the endpoint-window condition leaves a gap of at least
$3s_0/4$ in those windows.  Away from
these windows, the averaged mesh bound controls the sum of collision
probabilities over all intervals.  The maximum bound controls the same
sum where a bridge's time density is large near an outer endpoint; these
two uses are separated in Proposition~\ref{bridge:averaged-hit}.

\begin{lemma}[Free cutoff bounds]\label{path:regularity}
For every fixed $q>0$ and fixed $t>0$, a free Brownian path fails
regularity with probability at most $C_{q,t}D^{-q}$.
The bound is uniform in its initial point, $K$, and the allowed
$\alpha$-window.  It also holds for a Euclidean bridge of duration
$2t$ with bounded endpoint displacement, or for two independent
Euclidean bridge legs of duration $t$ joined at a prescribed middle
point, when the endpoints and middle point lie in a fixed bounded
region.  The bound and the required lower threshold on $D$ may depend on
that region and on $t$; the numerical regularity threshold $C_{\rm reg}$
is fixed once for all these laws.

If all three prescribed points lie within distance $3$ of a fixed
site $v$, the probability that the two-leg path leaves the ball of
radius $10$ about $v$ is at most $Ce^{-c/t}$, with numerical $c,C>0$.
\end{lemma}

\begin{proof}
For Brownian motion, the variables $A_j$ use independent increments.
Their tails are uniformly Gaussian, since $\tau/r^2\in[1/2,1]$.
Thus the maximum condition fails with probability at most
$C m_t e^{-cD^{2\gamma}}$.  Put $U_j=(1+A_j)^3$.
Their moments of every fixed order are uniformly bounded.  In the
expansion of the $2k$th moment of
$\sum_j(U_j-\E U_j)$, a nonzero term uses each of its indices at
least twice; there are at most $C_k m_t^k$ such terms.  Consequently
\[
 \E\left|\frac1{2m_t}\sum_j(U_j-\E U_j)\right|^{2k}
 \le C_k m_t^{-k}.
\]
Choose $C_{\rm reg}$ greater than a fixed multiple of
$1+\sup\E U_j$.  Since $m_t$ is comparable to $D^2$ at fixed $t$,
Markov's inequality proves every desired polynomial bound for the
average condition.  On an endpoint window the maximal Gaussian bound
has exponent a constant times
\[
 s_0^2/h_0=b^2D^{1/3},
\]
which proves the remaining free estimate.

For a bridge leg of duration $t$, write its centered part as
$W(u)-(u/t)W(t)$, using a Brownian motion $W$, and add its fixed
endpoint line.  The extra increment on a mesh interval, divided by
$r$, is at most
\[
 \frac{r}{t}\bigl(|W(t)|+C\bigr).
\]
Outside an event of arbitrarily small polynomial probability this
tends to zero, uniformly for bounded endpoint displacement.  The
inequality $(x+y)^3\le4x^3+4y^3$ transfers the averaged cutoff bound
after the numerical choice of $C_{\rm reg}$ is enlarged once.
On an endpoint window the extra displacement is
$h_0(|W(t)|+C)/t=o(s_0)$ with the same probability control.
These estimates prove the assertion for two legs, using their two
Brownian representations and a union bound.  A duration-$2t$ bridge
uses the identical representation with $2t$ in place of $t$.

Finally the endpoint lines of the two legs remain in the radius-$3$
ball.  Leaving the radius-$10$ ball requires a centered fluctuation
of size at least seven on one leg.  The Brownian representation and
the maximal Gaussian bound give the stated $Ce^{-c/t}$ estimate.
\end{proof}

\subsection{Many paths visiting a prescribed set}

For a set $A$ of lattice sites let $J_A(\omega)$ be the number of
sites $v\in A$ whose radius-$12$ ball the path visits.  Let $\mathcal D$
be a reversal-invariant single-path event and put
$\varepsilon_{\mathcal D}=\sup_y\PP^y_{\rm free}(\mathcal D)$.
All bounds below are uniform for $t\le1$ and sufficiently large $K$.

\begin{lemma}[A free visit bound]\label{path:free-visits}
For every $\lambda>0$,
\begin{equation}\label{path:free-visit-bound}
 \E^y_{\rm free}\bigl[e^{\lambda J_A\ind_{\mathcal D}}-1\bigr]
 \le \zeta e^{-d(y,A)},\qquad
 \zeta=C_\lambda\lambda\varepsilon_{\mathcal D}^{1/3}.
\end{equation}
Here $d(y,A)$ is the distance from $y$ to the set of site centers.
The constants $C_\lambda$ are bounded when $\lambda$ ranges over a
bounded interval.
\end{lemma}

\begin{proof}
Mark the lifted Brownian path at its initial point and at each first
exit from the unit ball about its last mark.  If $Q$ is the number
of completed exits by time $2t$, then $J_A\le C(Q+1)$: between
successive marks the path remains in a unit ball, so it can visit
radius-$12$ balls about only a bounded number of lattice sites.
Projection to the torus cannot increase this bound.

The exit durations are iid copies of a positive random variable
$\sigma$.  For $s>0$,
\[
 \PP(Q\ge k)\le e^{2s}(\E e^{-s\sigma})^k.
\]
Since $\E e^{-s\sigma}\to0$ as $s\to\infty$, $Q$ has every fixed
exponential moment, uniformly for $t\le1$.  Moreover $J_A>0$
requires a lifted displacement at least $(d(y,A)-12)_+$, and hence
has a Gaussian tail in that distance.  Using
$e^{\lambda J_A}-1\le\lambda J_A e^{\lambda J_A}$ and H\"older
with three equal exponents gives
\[
 \E[(e^{\lambda J_A}-1)\ind_{\mathcal D}\ind_{\{J_A>0\}}]
 \le \lambda\bigl(\E[J_A^3e^{3\lambda J_A}]\bigr)^{1/3}
       \varepsilon_{\mathcal D}^{1/3}
       \PP(J_A>0)^{1/3}.
\]
The exponential moment and Gaussian travel bound imply
Equation~\eqref{path:free-visit-bound}, after changing its constant.
\end{proof}

\begin{lemma}[Interacting visit moments]\label{path:visit-moments}
If $e^{Ct\alpha}\zeta\le1$, then
\begin{equation}\label{path:visit-mgf}
 \E_{\cR}\exp\left(\lambda\sum_i
             J_A(\omega_i)\ind_{\mathcal D}(\omega_i)\right)
 \le \exp\bigl(Ce^{Ct\alpha}\zeta D|A|\bigr).
\end{equation}
\end{lemma}

\begin{proof}
Set $h(\omega)=e^{\lambda J_A(\omega)\ind_{\mathcal D}(\omega)}-1$.
Expand $\prod_i(1+h(\omega_i))$ over subsets of labels and apply
Lemma~\ref{path:deletion} to each term.  Lemma~\ref{path:free-visits}
then bounds the result by
\[
 \int\mu(\dd Y)\prod_i
      [1+e^{Ct\alpha}\zeta e^{-d(Y_i,A)}]
 \le \int e^{e^{Ct\alpha}\zeta f_A}\dd\mu,
 \qquad f_A=\sum_i e^{-d(Y_i,A)}.
\]
Equation~\eqref{sp:count-mgf} applies with the indicated exponent
at most one, and proves Equation~\eqref{path:visit-mgf}.
\end{proof}

\subsection{Cells with a supply of movable paths}

Let $d_i(X)=\min_{j\ne i}d(X_i,X_j)$, taking the minimum to be
$+\infty$ when $N=1$.  Define the endpoint close count
\[
 C_v(X)=\#\{i:d(X_i,v)\le5,\ d_i(X)\le s_0\}.
\]
We call the following conditions a \emph{certificate at $v$}:
\begin{enumerate}[label=\textup{(\roman*)}]
 \item $n_v(Y)\ge\ell D$, and $C_v(Y),C_v(Z)<\theta D$;
 \item at most $H_1D$ trajectories visit the radius-$12$ ball about
 $v$, every visiting trajectory is regular, and fewer than $\theta D$
 of the visiting trajectories have lifted diameter greater than one.
\end{enumerate}
The certificate is an event of the full path sample.  Later it will
be used to estimate the probability of suitable exposed data; it
will not be added to the conditioning on unexposed trajectories.

\begin{proposition}[Joint certificate supply]\label{path:certificates}
Fix any $p_1>0$.  One can choose $H_1$ large, then $t>0$ sufficiently
small, and then $D_0$ sufficiently large, such that
\begin{equation}\label{path:certificate-bound}
 \cR\{v\text{ has no certificate for every }v\in A\}
 \le p_1^{|A|}
\end{equation}
for every set $A$ of distinct sites, every $D\ge D_0$, and sufficiently
large $K$.  The bound is uniform over the chosen ground function and
the allowed $\alpha$-window.  After $H_1$ is fixed, $t$ may be reduced
further before increasing $D_0$.
\end{proposition}

\begin{proof}
We bound simultaneous failures of each individual test first.
The low-count estimate in Proposition~\ref{sp:low-counts} and the
close-count estimate in Equation~\eqref{sp:close-tail} apply at both
endpoints by Lemma~\ref{path:slab}.  Their per-site parameters tend
to zero as $D$ increases.

For excessive visiting paths use $\mathcal D$ equal to all paths.
Choose a small fixed $\lambda>0$ so that
$e^{Ct\alpha}C_\lambda\lambda\le1$ for every $t\le1$ and the
allowed $\alpha$.  If every site in $A$ has more than $H_1D$
visitors, then $\sum_iJ_A(\omega_i)>H_1D|A|$.
Equation~\eqref{path:visit-mgf} and exponential Markov inequality give
\[
 \cR(\text{this event})
 \le\exp\{-[\lambda H_1-Ce^{Ct\alpha}C_\lambda\lambda]D|A|\}.
\]
Take $H_1$ large enough to make the bracket positive uniformly.

For visitors of diameter greater than one take $\lambda=1$ and
$\mathcal D=\{\diam\widetilde\omega>1\}$.
Its free probability is at most $Ce^{-c/t}$.
Choose $t$ small enough that the right-hand coefficient in
Equation~\eqref{path:visit-mgf} is at most $\theta/2$.
If every site in $A$ has at least $\theta D$ such visitors, their
sum of visits is at least $\theta D|A|$, so the probability is at
most $e^{-\theta D|A|/2}$.

Finally let $\mathcal D$ be failure of regularity.  For any fixed
large $\lambda$, Lemma~\ref{path:regularity} gives
$D\zeta\to0$.  If every site in $A$ has a nonregular visitor, the
sum of visits by such paths is at least $|A|$.  Thus
Equation~\eqref{path:visit-mgf} bounds this probability by
$\exp[-(\lambda-o(1))|A|]$.  Choosing $\lambda$ first and then
$D$ makes its per-site parameter as small as desired.  The same
increase of $D$ makes the preceding five individual failure bounds
arbitrarily small.

There are six failure types.  If every site of $A$ fails some test,
assign one type to each failing site.  In each of the at most
$6^{|A|}$ assignments, one type occurs on a subset of size at least
$|A|/6$.  If each single-type bound has parameter at most
$\varepsilon$, the union has probability at most
$(6\varepsilon^{1/6})^{|A|}$.  Choose
$\varepsilon\le(p_1/6)^6$.  This proves
Equation~\eqref{path:certificate-bound}.  The free estimates remain
valid when the already chosen $t$ is reduced and $D_0$ is then
increased, which proves the final assertion.
\end{proof}

We have obtained an unconditional joint bound for the cells where the
full sample is unsuitable.  The next section estimates a trial path
against a deterministic collection of such trajectories.  We will then
expose most paths in a way that preserves an exact conditional product
representation for the paths still to be changed.
